\section{The actual narrow common-sum domain is not preserved by the raw matrix}
\label{sec:smaller-sum-noninvariance}

MA-MO-005 concerns a coordinate core in the full smooth product.
The following result instead uses a genuine vector of Meyer's smaller
common-sum space. It determines its intersection with the narrow
Hilbert product and exhibits a literal raw output outside that product.
The operator is the original approximate projection; it is not replaced
by an orthogonal projection or an exact idempotent.

Fix \(0<a<1/2\), and put
\[
 \begin{aligned}
 H_{\Sigma,a}&=L^2\bigl(\mathbb R,
                         (e^{2ax}+e^{-2ax})\,dx\bigr),\\
 \mathcal S_>&=\{y:e^{bx}y\in\mathscr S(\mathbb R)
                                      \text{ for every }b>1/2\},\\
 \mathcal S_<&=\{y:e^{bx}y\in\mathscr S(\mathbb R)
                                      \text{ for every }b<-1/2\},\\
 \mathscr S_{\rm sum}
   &=\{(y_0,y_1)\in\mathcal S_>\oplus\mathcal S_<:\\
   &\hspace{2em}y_0-y_1=c_+e^{x/2}-c_-e^{-x/2}
                     \text{ for some }c_+,c_-\in\mathbb C\}.
 \end{aligned}
 \tag{MA-MO.78}
\]
These are the centered rational trivial-compact-isotype spaces of
Meyer (2005), \S\S5.2--5.3. The equation for the difference can also
be imposed in distributions on any locally square-integrable pair.

\paragraph{MA-MO-007 (exact intersection and original raw noninvariance).}
The pairs in \(\mathcal H_a\oplus\mathcal H_{-a}\) satisfying the
polar-difference equation in (MA-MO.78) are exactly
\[
 \{(y,y):y\in H_{\Sigma,a}\},\qquad
 \|(y,y)\|_{\mathcal H_a\oplus\mathcal H_{-a}}^2
                         =\|y\|_{H_{\Sigma,a}}^2.
 \tag{MA-MO.79}
\]
At the original smooth level,
\[
 \mathscr S_{\rm sum}\cap(\mathcal H_a\oplus\mathcal H_{-a})
                          =\{(y,y):y\in\mathcal E\}.
 \tag{MA-MO.80}
\]
Its Hilbert closure is (MA-MO.79).

Already for the source place set \(\{\infty\}\), the actual raw
operator \(P_+-P_-\) fails to preserve this narrow domain. There is
an even \(f_*\in C_c^\infty(\mathbb R\setminus\{0\})\) with
\(f_*(0)=\int f_*=0\), whose actual centered theta image
\(y_*\in\mathcal E\) satisfies
\[
 i_+f_*=(y_*,y_*),\qquad
 \bigl((P_+-P_-)(y_*,y_*)\bigr)_0\notin\mathcal H_a.
 \tag{MA-MO.81}
\]
The source compositions in this formula are all defined on their
original smooth domains. Consequently no Hilbert-valued operator
whose domain contains the entire original smooth diagonal can agree
with this raw matrix there. This is a failure of invariance of the
specified domain, not an assertion of nonclosability of an already
Hilbert-valued restriction.

\begin{proof}
\emph{The exact narrow intersection.}
If \(y_0\in\mathcal H_a\) and \(y_1\in\mathcal H_{-a}\), then
both \(e^{-ax}y_0\) and \(e^{-ax}y_1\) belong to
\(L^2((0,\infty))\): for the first use
\(e^{-ax}y_0=e^{-2ax}(e^{ax}y_0)\). Thus on this half-line
their difference belongs to \(L^2\). In the polar-difference
equation it equals
\(c_+e^{(1/2-a)x}-c_-e^{-(1/2+a)x}\). The second term
is square integrable, while the first is not unless \(c_+=0\).
On the negative half-line both \(e^{ax}y_j\) are square
integrable, since \(e^{ax}y_1=e^{2ax}(e^{-ax}y_1)\) there.
With \(c_+=0\), the remaining term
\(-c_-e^{(a-1/2)x}\) is not square integrable at the left
end unless \(c_-=0\). Equality in distributions gives the
same almost-everywhere equality, so this argument also covers
the stated distributional formulation. It proves (MA-MO.79).

The intersection \(\mathcal S_>\cap\mathcal S_<\) is exactly
\(\mathcal E\). To check the nontrivial inclusion, any desired
real exponential weight lies strictly between one permitted
left weight and one permitted right weight. On each half-line
its weighted derivatives are bounded by a decaying exponential
times the corresponding Schwartz estimates; a compact partition
joins the two bounds. This holds for every Schwartz seminorm.
Thus (MA-MO.80) follows from (MA-MO.79). The smooth compact
functions belong to \(\mathcal E\) and are dense in
\(H_{\Sigma,a}\), proving the closure assertion. The norm in
(MA-MO.79) is exactly the inherited norm; no equivalent
one-sided norm is substituted. It is also equivalent to the
usual weight \(e^{2a|x|}\), since the sum of the two
exponentials lies between that weight and twice that weight.

\emph{The actual archimedean source maps.}
Use self-dual additive Haar measure on \(\mathbb R\), and let
\(\mathfrak F\) have kernel \(e^{-2\pi it\xi}\). Work in the
even sector, where \(\mathfrak F^*=\mathfrak F\) and
\(\mathfrak F^2=1\). This is the source \(S=\{\infty\}\)
local model: the rational \(S\)-units are \(\{1,-1\}\).
Write the original cutoff in additive coordinates as
\[
 v(t)=\phi(\log|t|)\quad(t\ne0),\qquad v(0)=1,
                         \qquad q(t)=1-v(t).
 \tag{MA-MO.82}
\]
Because \(\phi\) equals one sufficiently far left and zero
sufficiently far right, \(v\) is an even compact smooth
function, equal to one near zero. Thus \(q=0\) near zero
and multiplication by \(q\) preserves Schwartz space.

The centered actual theta map is
\[
 \Theta f(x)=e^{x/2}\sum_{n\in\mathbb Z\setminus\{0\}}f(ne^x).
 \tag{MA-MO.83}
\]
Poisson summation gives the exact source identity
\[
 R\Theta\mathfrak F f
      =\Theta f+f(0)e^{x/2}
                     -\left(\int_{\mathbb R}f\right)e^{-x/2},
                          \qquad Ry(x)=y(-x).
 \tag{MA-MO.84}
\]
Hence \(i_+f=(\Theta f,R\Theta\mathfrak F^*f)\) is diagonal
when \(f(0)=\int f=0\). From the defining source maps,
not from an assumed idempotence, one has
\[
 p_+i_+f=qf+\mathfrak F(q\mathfrak F^*f),
                \qquad P_+=i_+p_+,\qquad
 P_-(y,y)=(y,y).
 \tag{MA-MO.85}
\]
The formula for \(p_+i_+\) is an approximate-section formula.
It is not simplified to the identity. Both summands on its
right are actual even Schwartz functions.

\emph{A genuine diagonal input with a polar output.}
Choose nonnegative \(\eta_L,\eta_R\in C_c^\infty(\mathbb R)\)
of integral one, supported respectively where \(\phi=1\)
and where \(\phi=0\). Set
\[
 \begin{aligned}
 F_*(x)&=e^{-x/2}(\eta_L(x)-\eta_R(x)),\\
 f_*(t)&=\frac{F_*(\log|t|)}{2\sqrt{|t|}}\\
       &=\frac{\eta_L(\log|t|)-\eta_R(\log|t|)}{2|t|}
                                     \quad(t\ne0),\qquad f_*(0)=0.
 \end{aligned}
 \tag{MA-MO.86}
\]
This is even, compactly supported away from zero, and smooth.
Changing variables \(t=e^x\) on the positive half-line shows
\(\int f_*=0\) and \(\int vf_*=1\). Both original theta
polar conditions therefore hold. The right tail of
\(y_*=\Theta f_*\) is rapidly decreasing by the actual sum,
and the left tail is rapidly decreasing by (MA-MO.84) and
Schwartz estimates for \(\mathfrak F f_*\). The same
estimates apply to all derivatives, so \(y_*\in\mathcal E\).
In particular \((y_*,y_*)\) lies in the exact domain
(MA-MO.80), not merely in the full smooth product.

Let \(h=p_+i_+f_*\), with its literal value from (MA-MO.85).
Its additive integral is
\[
 \int_{\mathbb R}h
   =\int qf_*+(q\mathfrak F^*f_*)(0)
   =-1.
 \tag{MA-MO.87}
\]
Here the second term is zero because \(q=0\) near zero;
the first is \(\int f_*-\int vf_*=-1\). Fourier inversion
at zero justifies the first equality for these Schwartz
functions. In particular the image has a genuine nonzero
source polar term.

For any even Schwartz \(h\), the full Poisson formula gives
\(\Theta h(x)=(\int h)e^{-x/2}-h(0)e^{x/2}\) plus a
remainder rapidly decreasing with all derivatives as
\(x\to-\infty\). Apply this to the present \(h\). The
first component of \((P_+-P_-)(y_*,y_*)\) is
\(\Theta h-y_*\), whose left asymptotic is
\(-e^{-x/2}-h(0)e^{x/2}\) plus a rapid remainder. Its
leading coefficient is nonzero, so it does not belong to
\(\mathcal H_a\) for \(0<a<1/2\). This proves (MA-MO.81)
on the actual source domain and completes the proof.
\end{proof}

\paragraph{Compact smoothing and the value-zero test condition.}
The same example also gives an exact scope statement for smoothing.
If \(g\in C_c^\infty(\mathbb R)\), compact convolution preserves
the displayed asymptotics and acts on each polar exponential by
its Mellin value. Consequently
\[
 \bigl(C_g(P_+-P_-)(y_*,y_*)\bigr)_0
   =-\mathcal Mg(1/2)e^{-x/2}
     -h(0)\mathcal Mg(-1/2)e^{x/2}
       +\text{a rapid left remainder}.
 \tag{MA-MO.88}
\]
If \(\mathcal Mg(1/2)\ne0\), this raw smoothed output is
still outside \(\mathcal H_a\). The condition \(g(0)=0\)
does not remove that term: a nonzero nonnegative test supported
in \((1,2)\) satisfies \(g(0)=0\) but has positive
\(\mathcal Mg(1/2)\). No conflict with MA-RS is implied,
since its continued operators and explicitly subtracted polar
terms have different, specified domains.

\paragraph{Further pole-free restriction and source boundary.}
The result does not decide the norm of the raw action after
restricting its inputs further so that the output also has
both polar coefficients zero. Nor does it identify the closure
of such a further domain, or infer nonclosability from an
output that is not in the proposed Hilbert target. These are
separate questions. In particular MA-MO-007 is a genuine
common-sum domain test, not a reuse of the coordinate-axis
witness from MA-MO-005.

The exact primary locators are Meyer (2005),
\texttt{sec:summation}, \texttt{lem:Hilm\_p\_embed}, the
display for \(p_+^S i_+\), the defining matrices \(P_+^S,P_-\),
\texttt{eq:Poisson\_summation}, and \texttt{lem:Poisson\_sumpm};
author-manuscript pp.~39--41, supplied \texttt{main.tex}
lines~3808--4043. The cutoff orientation is also explicit
at lines~1403--1407 and~1703--1709. The sign is the one
derived from the defining matrices in MA-MEY-003.
Meyer's Lemma~5.6 concerns centered weights \(a>1/2\),
where these source polar tails have the allowed weighted
behavior, and is unaffected by this narrow-domain result.

The two-dimensional singular representation is part of the
actual source exact-sequence geometry. Its polar generators
are absent from the inherited narrow Hilbert intersection,
as (MA-MO.79) proves. Adding them by a new norm or removing
them by a chosen finite-rank correction changes that analytic
realization and requires its own map. Neither the finite-length
spectral \(K_0\) class nor the inherited cote~1 nuclear-trace
locators, scans~079--080 and~111--112 through
ALPOGE\_BRIDGES BR-TR-01/02, supply that analytic change.
No A, M, or other zero-location premise is used here.
